%\ifodd\value{page}\hbox{}\newpage\fi
\chapter*{Ṛgveda 10.125 — Devī Sūkta, Mantra 6}
\addcontentsline{toc}{chapter}{Devī Sūkta 10.125.6}

\begin{sloka}
{\sanskritfont
\begin{tabular}{c}
अहं सुवे पितरमस्य मूर्धन्मम योनिरप्स्वन्तः समुद्रे ।\\
ततो वि तिष्ठे भुवनानु विश्वोतामूं द्यां वर्ष्मणोपस्पृशामि ॥६॥
\end{tabular}

\vspace{1.5em}

{\middlesize \iastfont
ahaṁ suve pitaram asya mūrdhan mama yonir apsv antaḥ samudre |\\
tato vi tiṣṭhe bhuvanānu viśvotāmūṁ dyāṁ varṣmaṇopaspṛśāmi ||6||}

\vspace{1em}

{\middlesize
\anvaya{%
अहं अस्य पितरम् मूर्धन् सुवे ।
मम योनिः अप्सु अन्तः समुद्रे ।
ततः अहं भुवनानि अनु विश्वम् वितिष्ठे ।
उत अमूम् द्याम् वर्ष्मणा उपस्पृशामि ।
}
}

\middlesize \iastfont \itmean{%
«Io genero il padre di questo mondo, al suo vertice;  
il mio grembo è nelle acque, nel profondo dell’oceano.  
Da lì io mi estendo attraverso tutti gli esseri,  
e tocco quel cielo lassù con la mia grandezza.»%
}

}
\end{sloka}

\wordmeaning{%
\wordline{अहं सुवे पितरम् अस्य मूर्धन्}
{ahaṁ suve pitaram asya mūrdhan}
{ io genero (\textit{suve}) il padre (\textit{pitaram}) di questo mondo, al suo capo o vertice (\textit{mūrdhan}); }%

\wordline{मम योनिः अप्सु अन्तः समुद्रे}
{mama yoniḥ apsu antaḥ samudre}
{ il mio grembo (\textit{yoni}) è nelle acque (\textit{apsu}), nel profondo (\textit{antaḥ}) dell’oceano (\textit{samudra}); }%

\wordline{ततः वि तिष्ठे भुवनानि अनु विश्वम्}
{tataḥ vi tiṣṭhe bhuvanāni anu viśvam}
{ da lì (\textit{tataḥ}) io mi dispiego (\textit{vi-tiṣṭhe}) attraverso (\textit{anu}) tutti i mondi e gli esseri (\textit{bhuvana}); }%

\wordline{उत अमूम् द्याम् वर्ष्मणा उपस्पृशामि}
{uta amūṁ dyām varṣmaṇā upaspṛśāmi}
{ e anche (\textit{uta}) tocco (\textit{upaspṛśāmi}) quel cielo (\textit{amūṁ dyām}) con la mia grandezza, estensione (\textit{varṣman}); }%
}

Nel sesto mantra la Dea articola una cosmologia completa. Lei genera il “padre” del mondo, collocandosi così come principio anteriore a ogni paternità: non è effetto, ma sorgente. Il suo grembo è situato nelle acque primordiali, luogo indifferenziato e generativo da cui emerge l’ordine cosmico. Da questo centro nascosto, la Dea si dispiega attraverso tutti gli esseri e tutti i mondi. La sua presenza non è confinata né localizzata: è una diffusione continua che attraversa il cosmo dall’interno. L’atto finale di “toccare il cielo” con la propria grandezza indica che la sua estensione non ha limite superiore: lei unisce il profondo e l’alto, l’oceano e il cielo. Questo mantra conclude il Devī Sūkta con una visione di totalità: la Dea come grembo, origine, diffusione e trascendenza. Ciò che genera, ciò che sostiene e ciò che supera coincidono in un’unica potenza, che parla in prima persona come fondamento ultimo del reale.

\textit{Yoni} indica il grembo, la sorgente generativa, ma nel contesto vedico il suo significato trascende ogni riferimento puramente biologico: è lo spazio dell’indifferenziato fecondo, il luogo in cui le forme non sono ancora separate e da cui ogni forma può emergere. Collocare la \textit{yoni} “nelle acque, nel profondo dell’oceano” significa situare l’origine del cosmo in una profondità primordiale, anteriore a ogni distinzione tra cielo e terra, soggetto e oggetto, maschile e femminile. La Dea non nasce dal mondo: è il mondo che nasce dal suo grembo. In questo senso, \textit{yoni} precede persino la paternità, poiché ella genera “il padre di questo mondo”, mostrando che ogni principio ordinatore deriva da una matrice più antica e più vasta. \textit{Yoni} è una origine che non si esaurisce nell’atto di creare, ma rimane attiva come potenza che permea e mantiene l’intero cosmo.



